\documentclass{article}
\usepackage[UTF8]{ctex}
\usepackage{amsmath,amssymb,amsfonts}
\usepackage{enumitem}
\usepackage{geometry}
\geometry{a4paper, margin=1in}
% === 错题本排版优化（不参与题目标识） ===
\setlength{\parindent}{0pt}
\setlength{\parskip}{0.35em plus 0.10em minus 0.08em}
\setlength{\emergencystretch}{3em}
\setlist{itemsep=0.25em, topsep=0.40em, parsep=0pt}
\allowdisplaybreaks[2]
% === 错题本排版优化结束 ===
\begin{document}
\section*{华中师范大学2024年数学分析真题}
\section*{138．华中师范大学2024年数学分析试卷}
\begin{itemize}

\subsection{题目1}
  \item[1．] 计算
\[
\lim _{x \rightarrow+\infty}\left[\left(x^{3}-x^{2}+\frac{x}{2}\right) e^{\frac{1}{x}}-\sqrt{1+x^{6}}\right] .
\]

\subsection{题目2}
  \item[2．] 将 \(f(x)=\sin x, x \in[0, \pi]\) 展开成余弦级数．
\end{itemize}

\subsection{题目3}
3．计算 \(\iint_{D} e^{\frac{y}{z+v}} \mathrm{~d} x \mathrm{~d} y\) ，其中 \(D=\{(x, y) \mid x+y \leq 1, x \geq 0, y \geq 0\}\) ．

\subsection{题目4}
4．计算

\[
I=\int_{\Gamma} \frac{(x+y) \mathrm{d} x-(x-y) \mathrm{d} y}{x^{2}+y^{2}}
\]

其中 \(\Gamma\) 是由方程 \(x^{2}+y^{2}=a^{2}(a>0)\) 确定的曲线，方向取逆时针．

\subsection{题目5}
5．计算

\[
I=\iint_{S} x \mathrm{~d} x \mathrm{~d} y+y \mathrm{~d} z \mathrm{~d} x+z \mathrm{~d} x \mathrm{~d} y
\]

其中 \(S\) 是由方程 \(x^{2}+y^{2}+z^{2}=a^{2}(a>0)\) 确定的曲面，方向取外侧．

二．（15 分）设 \(\lim _{n \rightarrow \infty} a_{n}=a, \lim _{n \rightarrow \infty} b_{n}=b\) ，记

\[
c_{n}=\frac{a_{1} b_{n}+a_{2} b_{n-1}+\cdots+a_{n} b_{1}}{n} .
\]

请用数列极限的定义证明 \(\lim _{n \rightarrow \infty} c_{n}=a b\) 。

三．（10 分）设 \(f(x)\) 在 \((-\infty,+\infty)\) 上二阶可微，记 \(M_{i}=\sup _{-\infty<x<+\infty}\left|f^{(i)}(x)\right|<+\infty(i=0,1,2)\) ．证明：

\[
M_{1}^{2} \leq 2 M_{0} M_{2}
\]

四．（15 分）设 \((a, b)\) 为有界区间，证明：函数 \(f(x)\) 在 \((a, b)\) 上一致连续的充要条件是对于 \((a, b)\) 上的任意 Cauchy 列 \(\left\{x_{n}\right\}\) ，函数值数列 \(\left\{f\left(x_{n}\right)\right\}\) 也是 Cauchy 列。

五．（15 分）设 \(f(x)\) 为 \(\mathbb{R}\) 上的连续函数，记 \(\varphi(x)=f(x) \int_{0}^{x} f(t) \mathrm{d} t, x \in(-\infty,+\infty)\) ，若 \(\varphi(x)\) 在\((-\infty,+\infty)\) 上单调递减，证明：\(f(x) \equiv 0\) ．

六．（10 分）设 \(\left\{x_{n}\right\}\) 为 \((0,1)\) 中的数列，满足当 \(i \neq j\) 时，\(x_{i} \neq x_{j}\) ，试讨论函数 \(f(x)=\sum_{n=1}^{\infty} \frac{\operatorname{sgn}\left(x-x_{n}\right)}{2^{n}}\) 在\((0,1)\) 上的连续性，其中 \(\operatorname{sgn}(x)\) 表示符号函数．

七．（15 分）含参量反常积分 \(I(r)=\int_{0}^{+\infty} r e^{-r x} \mathrm{~d} x\) ．

\begin{itemize}
  \item[（1）] 证明：\(I(r)\) 在区间 \([a, b]\) 上关于 \(r\) 一致收敛，其中 \(0<a<b\) ．
  \item[（2）] 证明：\(I(r)\) 在区间 \([0, b]\) 上关于 \(r\) 不一致收敛，其中 \(0<b\) ．
\end{itemize}

八．（20分）设级数 \(\sum_{n=1}^{\infty} a_{n}\) 收敛，级数 \(\sum_{n=1}^{\infty}\left(b_{n+1}-b_{n}\right)\) 绝对收敛。

\begin{itemize}
  \item[（1）] 叙述阿贝尔变换，即级数的分部求和公式。
  \item[（2）] 证明： \(\lim _{n \rightarrow \infty} b_{n}\) 存在．
  \item[（3）] 证明：\(\sum_{n=1}^{\infty} a_{n} b_{n}\) 收敛．
\end{itemize}
\end{document}
